\documentclass{article}
\usepackage[utf8]{inputenc}
\usepackage[a4paper, margin=1in]{geometry}
\usepackage{listings}
\usepackage{xcolor}

\definecolor{codegreen}{rgb}{0,0.6,0}
\definecolor{codegray}{rgb}{0.5,0.5,0.5}
\definecolor{codepurple}{rgb}{0.58,0,0.82}
\definecolor{backcolour}{rgb}{0.95,0.95,0.92}

\lstdefinestyle{mystyle}{
    backgroundcolor=\color{backcolour},
    commentstyle=\color{codegreen},
    keywordstyle=\color{magenta},
    numberstyle=\tiny\color{codegray},
    stringstyle=\color{codepurple},
    basicstyle=\ttfamily\footnotesize,
    breakatwhitespace=false,
    breaklines=true,
    captionpos=b,
    keepspaces=true,
    numbers=left,
    numbersep=5pt,
    showspaces=false,
    showstringspaces=false,
    showtabs=false,
    tabsize=2
}
\lstset{style=mystyle}

\title{\textbf{Lab 5.5: Remaining 10 SQL Queries}}
\author{Name: [Your Name Here] \\ ID: [Your ID Here]}
\date{\today}

\begin{document}

\maketitle

\section*{Queries 31 to 40}

\begin{lstlisting}[language=SQL, caption=Lab 5 Additional Queries]
USE MedicalShopDB;

-- 31. DDL: Create Prescription Table
CREATE TABLE IF NOT EXISTS Prescription (
    prescription_id INT PRIMARY KEY AUTO_INCREMENT,
    customer_id     INT NOT NULL,
    medicine_id     INT NOT NULL,
    prescribed_on   DATE NOT NULL,
    dosage          VARCHAR(100),
    FOREIGN KEY (customer_id) REFERENCES Customer(customer_id) ON DELETE CASCADE,
    FOREIGN KEY (medicine_id) REFERENCES Medicine(medicine_id) ON DELETE CASCADE
);

-- 32. DML: Insert into Prescription
INSERT INTO Prescription (customer_id, medicine_id, prescribed_on, dosage) VALUES
(1, 1, '2023-09-15', 'Twice a day after meals'),
(2, 2, '2023-10-01', 'Once at night'),
(3, 1, '2023-10-10', 'Thrice a day');

-- 33. DQL: INNER JOIN across three tables
SELECT
    p.prescription_id,
    c.name AS customer_name,
    m.type AS medicine_type,
    p.prescribed_on,
    p.dosage
FROM Prescription p
JOIN Customer c ON p.customer_id = c.customer_id
JOIN Medicine m ON p.medicine_id = m.medicine_id;

-- 34. DQL: WHERE with BETWEEN and LIKE
SELECT name, cp, sp, expiry
FROM Shop_Inventory
WHERE sp BETWEEN 10 AND 50
  AND name LIKE 'C%';

-- 35. DQL: GROUP BY with aggregate functions
SELECT
    m.type AS medicine_type,
    SUM(si.qty) AS total_qty,
    ROUND(AVG(si.sp), 2) AS avg_selling_price
FROM Shop_Inventory si
JOIN Medicine m ON si.medicine_id = m.medicine_id
GROUP BY m.type;

-- 36. DQL: HAVING filter after aggregation
SELECT
    payment_method,
    SUM(past_orders) AS total_orders
FROM Customer
GROUP BY payment_method
HAVING total_orders > 4;

-- 37. DQL: Subquery to find medicines stocked by wholeseller in bulk
SELECT name, qty, sp
FROM Shop_Inventory
WHERE medicine_id IN (
    SELECT medicine_id
    FROM Wholeseller_Inventory
    WHERE qty > 1000
);

-- 38. DML: Update selling price for low-margin items
UPDATE Shop_Inventory
SET sp = ROUND(sp * 1.10, 2)
WHERE (sp - cp) / cp < 0.50;

-- 39. View: Create Medicine Profit View
CREATE OR REPLACE VIEW Medicine_Profit_View AS
SELECT
    si.name,
    si.qty,
    si.cp,
    si.sp,
    ROUND((si.sp - si.cp), 2) AS unit_profit,
    ROUND((si.sp - si.cp) * si.qty, 2) AS total_profit
FROM Shop_Inventory si;

SELECT * FROM Medicine_Profit_View ORDER BY total_profit DESC;

-- 40. Trigger: Log when a bill status changes to Paid
CREATE TABLE IF NOT EXISTS Bill_Audit (
    audit_id INT PRIMARY KEY AUTO_INCREMENT,
    bill_id INT,
    customer_id INT,
    paid_on TIMESTAMP DEFAULT CURRENT_TIMESTAMP,
    amount DECIMAL(10,2)
);

DELIMITER //
CREATE TRIGGER after_bill_status_update
AFTER UPDATE ON Bills
FOR EACH ROW
BEGIN
    IF NEW.status = 'Paid' AND OLD.status <> 'Paid' THEN
        INSERT INTO Bill_Audit (bill_id, customer_id, amount)
        VALUES (NEW.bill_id, NEW.customer_id, NEW.total_amount);
    END IF;
END //
DELIMITER ;

UPDATE Bills SET status = 'Paid' WHERE bill_id = 2;
SELECT * FROM Bill_Audit;
\end{lstlisting}

\end{document}
